\begin{table}[t]\begin{threeparttable}
\centering
\caption{Standardized Effect Sizes}
\label{tab:sde}
\small
\begin{tabular}{lcccccc}
\toprule
  Outcome & $\hat{\beta}$ & SE & SD($Y$) & SDE & SE(SDE) & Classification \\
\midrule
\multicolumn{7}{l}{\textit{Panel A: Pooled}} \\[3pt]
  Log tax revenue & -0.0045 & 0.0006 & 0.868 & -0.053 & 0.008 & Moderate negative \\
  Log assessment & -0.0085 & 0.0007 & 0.767 & -0.112 & 0.010 & Moderate negative \\
  Forced sales & -0.343 & 0.097 & 19.6 & -0.178 & 0.051 & Large negative \\
  Grundskyld rate & 0.240 & 0.029 & 4.6 & 0.528 & 0.063 & Large positive \\[6pt]
\multicolumn{7}{l}{\textit{Panel B: Heterogeneous (by dose intensity)}} \\[3pt]
  Log tax (high-dose munis) & -0.0022 & 0.0021 & 1.013 & -0.013 & 0.012 & Small negative \\
  Log tax (low-dose munis) & -0.0037 & 0.0015 & 0.636 & -0.037 & 0.015 & Small negative \\
\bottomrule
\end{tabular}
\begin{tablenotes}
\item \textit{Notes:} \textbf{Country:} Denmark. \textbf{Research question:} Does the 2024 Boligskattereform's municipality-varying grundskyld rate reductions affect property tax revenue, assessed values, and forced sales across 98 Danish municipalities? \textbf{Policy mechanism:} The reform simultaneously reassessed all properties using 2022 valuations, slashed grundskyld rates by 32--88\%, and introduced a permanent tax-freeze loan for incumbents that resets upon sale, creating municipality-varying wedges in effective tax burdens. \textbf{Outcome definition:} Log total municipal property tax revenue (1000 DKK), log property assessment value (million DKK), annual count of forced property sales, and grundskyld rate (per mille of assessed land value). \textbf{Treatment:} Continuous; percentage change in municipality grundskyld rate from 2023 to 2024, ranging from $-87.5\%$ to $-32.3\%$. \textbf{Data:} Statistics Denmark open API (api.statbank.dk), tables EJDSK2, ESKAT, TVANG3; 98 municipalities, 2016--2025, 980 municipality-year observations. \textbf{Method:} Dose-response difference-in-differences with municipality and year fixed effects; standard errors clustered by municipality. \textbf{Sample:} All 98 Danish municipalities with complete grundskyld rate data 2016--2025; excludes regional and national aggregates. SDE $= \hat{\beta} \times \text{SD}(X) / \text{SD}(Y)$ where SD($X$) is the cross-municipality standard deviation of the dose and SD($Y$) is the pre-treatment standard deviation of the outcome. Classification refers to magnitude, not statistical significance: Large ($|$SDE$| > 0.15$), Moderate ($0.05$--$0.15$), Small ($0.005$--$0.05$), Null ($< 0.005$).
\end{tablenotes}
\end{threeparttable}\end{table}
